\documentclass[11pt,a4paper]{article}

\usepackage[utf8]{inputenc}
\usepackage[T1]{fontenc}
\usepackage[french]{babel}
\usepackage{amsmath, amssymb, amsthm}
\usepackage{geometry}
\usepackage{booktabs}
\usepackage{listings}
\usepackage{xcolor}
\usepackage{graphicx}

\geometry{margin=2.5cm}

\title{Corrigé - TD 5 : Probabilités}
\author{MT331 - INP Esisar Valence}
\date{Année 2025-2026}

\begin{document}

\maketitle

\section*{Exercice 1 - Couple Discret}

\subsection*{a) Déterminer les lois marginales du couple.}

\textbf{Raisonnement :} 
La probabilité marginale $\mathbb{P}(X=x)$ s'obtient en sommant les probabilités de la ligne correspondante à $x$.
La probabilité marginale $\mathbb{P}(Y=y)$ s'obtient en sommant les probabilités de la colonne correspondante à $y$.

\textbf{Reponse :}
Loi marginale de X :
\begin{itemize}
    \item $\mathbf{\mathbb{P}(X=0)} = 0.03 + 0.05 + 0.05 + 0.04 + 0.03 = \mathbf{0.20}$
    \item $\mathbf{\mathbb{P}(X=1)} = 0.07 + 0.08 + 0.10 + 0.05 + 0.05 = \mathbf{0.35}$
    \item $\mathbf{\mathbb{P}(X=2)} = 0.05 + 0.07 + 0.09 + 0.08 + 0.01 = \mathbf{0.30}$
    \item $\mathbf{\mathbb{P}(X=3)} = 0 + 0.05 + 0.06 + 0.03 + 0.01 = \mathbf{0.15}$
\end{itemize}

Loi marginale de Y :
\begin{itemize}
    \item $\mathbf{\mathbb{P}(Y=0)} = 0.03 + 0.07 + 0.05 + 0 = \mathbf{0.15}$
    \item $\mathbf{\mathbb{P}(Y=1)} = 0.05 + 0.08 + 0.07 + 0.05 = \mathbf{0.25}$
    \item $\mathbf{\mathbb{P}(Y=2)} = 0.05 + 0.10 + 0.09 + 0.06 = \mathbf{0.30}$
    \item $\mathbf{\mathbb{P}(Y=3)} = 0.04 + 0.05 + 0.08 + 0.03 = \mathbf{0.20}$
    \item $\mathbf{\mathbb{P}(Y=4)} = 0.03 + 0.05 + 0.01 + 0.01 = \mathbf{0.10}$
\end{itemize}

\subsection*{b) Calculer $\mathbb{E}(X)$, $\mathbb{V}(X)$, $\mathbb{E}(Y)$, $\mathbb{V}(Y)$ et $Cov(X;Y)$.}

\textbf{Raisonnement :} 
\textbf{Etape 1 : Espérance et Variance de X.}\\
$\mathbb{E}(X) = 0(0.20) + 1(0.35) + 2(0.30) + 3(0.15) = 1.4$.
$\mathbb{E}(X^2) = 0(0.20) + 1^2(0.35) + 2^2(0.30) + 3^2(0.15) = 2.9$.
$\mathbb{V}(X) = \mathbb{E}(X^2) - (\mathbb{E}(X))^2$.

\textbf{Etape 2 : Espérance et Variance de Y.}\\
$\mathbb{E}(Y) = 0(0.15) + 1(0.25) + 2(0.30) + 3(0.20) + 4(0.10) = 1.85$.
$\mathbb{E}(Y^2) = 0 + 1(0.25) + 4(0.30) + 9(0.20) + 16(0.10) = 4.85$.
$\mathbb{V}(Y) = \mathbb{E}(Y^2) - (\mathbb{E}(Y))^2$.

\textbf{Etape 3 : Calcul de la Covariance.}\\
On calcule d'abord $\mathbb{E}(XY) = \sum_{x,y} x \cdot y \cdot \mathbb{P}(X=x, Y=y)$. Seuls les cas où $x \neq 0$ et $y \neq 0$ comptent :
\begin{itemize}
    \item $x=1$ : $1(0.08) + 2(0.10) + 3(0.05) + 4(0.05) = 0.63$
    \item $x=2$ : $2(1(0.07) + 2(0.09) + 3(0.08) + 4(0.01)) = 2(0.53) = 1.06$
    \item $x=3$ : $3(1(0.05) + 2(0.06) + 3(0.03) + 4(0.01)) = 3(0.30) = 0.90$
\end{itemize}
$\mathbb{E}(XY) = 0.63 + 1.06 + 0.90 = 2.59$.
$Cov(X,Y) = \mathbb{E}(XY) - \mathbb{E}(X)\mathbb{E}(Y) = 2.59 - (1.4 \times 1.85)$.

\textbf{Reponse :}
$$ \mathbf{\mathbb{E}(X) = 1.4} \quad \text{et} \quad \mathbf{\mathbb{V}(X) = 0.94} $$
$$ \mathbf{\mathbb{E}(Y) = 1.85} \quad \text{et} \quad \mathbf{\mathbb{V}(Y) = 1.4275} $$
$$ \mathbf{Cov(X,Y) = 0} $$

\subsection*{c) Les variables aléatoires X et Y sont-elles indépendantes ?}

\textbf{Raisonnement :} 
Bien que la covariance soit nulle, cela ne suffit pas à prouver l'indépendance. On vérifie sur une case de la loi jointe. Prenons $X=3$ et $Y=0$.
$\mathbb{P}(X=3, Y=0) = 0$.
Pourtant, le produit des probabilités marginales donne : $\mathbb{P}(X=3) \times \mathbb{P}(Y=0) = 0.15 \times 0.15 = 0.0225$.
Puisque $0 \neq 0.0225$, la définition n'est pas vérifiée.

\textbf{Reponse :} 
\textbf{Les variables X et Y ne sont pas indépendantes.}

\subsection*{d) On pose $Z=X+2Y$. Déterminer l'espérance et la variance de Z.}

\textbf{Raisonnement :} 
Par la propriété de linéarité de l'espérance : $\mathbb{E}(Z) = \mathbb{E}(X) + 2\mathbb{E}(Y)$.
Par la formule de la variance pour une somme : $\mathbb{V}(Z) = \mathbb{V}(X) + 4\mathbb{V}(Y) + 4 Cov(X, Y)$.
Puisque $Cov(X, Y) = 0$, on remplace simplement par les valeurs trouvées en b).

\textbf{Reponse :}
$$ \mathbf{\mathbb{E}(Z) = 1.4 + 2(1.85) = 5.1} $$
$$ \mathbf{\mathbb{V}(Z) = 0.94 + 4(1.4275) + 0 = 6.65} $$

\subsection*{e) Déterminer la loi de Z.}

\textbf{Raisonnement :} 
$Z = x + 2y$. En balayant le tableau, les valeurs possibles vont de $0 + 0 = 0$ à $3 + 2(4) = 11$.
On additionne les probabilités des cases $(x,y)$ du tableau initial qui donnent la même valeur pour $x+2y$.

\textbf{Reponse :}
\begin{itemize}
    \item $\mathbf{\mathbb{P}(Z=0)} : (0,0) \implies \mathbf{0.03}$
    \item $\mathbf{\mathbb{P}(Z=1)} : (1,0) \implies \mathbf{0.07}$
    \item $\mathbf{\mathbb{P}(Z=2)} : (2,0), (0,1) \implies 0.05 + 0.05 = \mathbf{0.10}$
    \item $\mathbf{\mathbb{P}(Z=3)} : (3,0), (1,1) \implies 0 + 0.08 = \mathbf{0.08}$
    \item $\mathbf{\mathbb{P}(Z=4)} : (2,1), (0,2) \implies 0.07 + 0.05 = \mathbf{0.12}$
    \item $\mathbf{\mathbb{P}(Z=5)} : (3,1), (1,2) \implies 0.05 + 0.10 = \mathbf{0.15}$
    \item $\mathbf{\mathbb{P}(Z=6)} : (2,2), (0,3) \implies 0.09 + 0.04 = \mathbf{0.13}$
    \item $\mathbf{\mathbb{P}(Z=7)} : (3,2), (1,3) \implies 0.06 + 0.05 = \mathbf{0.11}$
    \item $\mathbf{\mathbb{P}(Z=8)} : (2,3), (0,4) \implies 0.08 + 0.03 = \mathbf{0.11}$
    \item $\mathbf{\mathbb{P}(Z=9)} : (3,3), (1,4) \implies 0.03 + 0.05 = \mathbf{0.08}$
    \item $\mathbf{\mathbb{P}(Z=10)} : (2,4) \implies \mathbf{0.01}$
    \item $\mathbf{\mathbb{P}(Z=11)} : (3,4) \implies \mathbf{0.01}$
\end{itemize}

\subsection*{f) Déterminer $Cov(X;Z)$.}

\textbf{Raisonnement :} 
On utilise la propriété de bilinéarité de la covariance :
$$ Cov(X, Z) = Cov(X, X+2Y) = Cov(X,X) + 2Cov(X,Y) $$
Or, la covariance d'une variable avec elle-même est sa variance ($Cov(X,X) = \mathbb{V}(X)$), et nous avons déjà démontré que $Cov(X,Y) = 0$.

\textbf{Reponse :}
$$ \mathbf{Cov(X, Z) = \mathbb{V}(X) = 0.94} $$

\hrulefill

\section*{Exercice 2 - Couple Discret}

\subsection*{a) Déterminer les lois marginales du couple et préciser si ces variables aléatoires sont indépendantes.}

\textbf{Raisonnement :} 
On somme les lignes pour obtenir la loi de X, et les colonnes pour obtenir la loi de Y.
Ensuite, pour l'indépendance, on vérifie si chaque case du tableau correspond au produit exact de la probabilité de sa ligne et de sa colonne.
Par exemple : $\mathbb{P}(X=1) \times \mathbb{P}(Y=1) = 0.40 \times 0.20 = 0.08 = \mathbb{P}(X=1,Y=1)$. C'est le cas pour toutes les valeurs.

\textbf{Reponse :}
Loi marginale de X :
\begin{itemize}
    \item $\mathbf{\mathbb{P}(X=1)} = 0.08 + 0.04 + 0.16 + 0.12 = \mathbf{0.40}$
    \item $\mathbf{\mathbb{P}(X=2)} = 0.04 + 0.02 + 0.08 + 0.06 = \mathbf{0.20}$
    \item $\mathbf{\mathbb{P}(X=3)} = 0.08 + 0.04 + 0.16 + 0.12 = \mathbf{0.40}$
\end{itemize}
Loi marginale de Y :
\begin{itemize}
    \item $\mathbf{\mathbb{P}(Y=1)} = 0.08 + 0.04 + 0.08 = \mathbf{0.20}$
    \item $\mathbf{\mathbb{P}(Y=2)} = 0.04 + 0.02 + 0.04 = \mathbf{0.10}$
    \item $\mathbf{\mathbb{P}(Y=3)} = 0.16 + 0.08 + 0.16 = \mathbf{0.40}$
    \item $\mathbf{\mathbb{P}(Y=4)} = 0.12 + 0.06 + 0.12 = \mathbf{0.30}$
\end{itemize}
\textbf{Les variables X et Y sont indépendantes.}

\subsection*{b) Calculer $Cov(X;Y)$.}

\textbf{Raisonnement :} 
D'après le cours, si deux variables aléatoires sont indépendantes, leur covariance est strictement nulle. Ayant prouvé l'indépendance à la question précédente, le calcul est direct.

\textbf{Reponse :}
$$ \mathbf{Cov(X, Y) = 0} $$

\subsection*{c) Déterminer la loi du couple $(U, V)$ où $U = \min(X;Y)$ et $V = \max(X;Y)$.}

\textbf{Raisonnement :} 
$X \in \{1,2,3\}$ et $Y \in \{1,2,3,4\}$. Par définition, on a toujours $U \le V$. Les valeurs possibles pour $U$ sont $\{1,2,3\}$ et pour $V$ sont $\{1,2,3,4\}$.
Pour construire les probabilités jointes $\mathbb{P}(U=u, V=v)$ :
\begin{itemize}
    \item Si $u = v$ : la probabilité est celle de la diagonale $\mathbb{P}(X=u, Y=u)$.
    \item Si $u < v$ : on additionne les deux cas symétriques, car soit X est le min, soit Y est le min.
    $\mathbb{P}(U=u, V=v) = \mathbb{P}(X=u, Y=v) + \mathbb{P}(X=v, Y=u)$.
\end{itemize}

\textbf{Reponse :}
\begin{itemize}
    \item $\mathbf{\mathbb{P}(U=1, V=1)} = \mathbb{P}(X=1,Y=1) = \mathbf{0.08}$
    \item $\mathbf{\mathbb{P}(U=1, V=2)} = \mathbb{P}(X=1,Y=2) + \mathbb{P}(X=2,Y=1) = \mathbf{0.08}$
    \item $\mathbf{\mathbb{P}(U=1, V=3)} = \mathbb{P}(X=1,Y=3) + \mathbb{P}(X=3,Y=1) = \mathbf{0.24}$
    \item $\mathbf{\mathbb{P}(U=1, V=4)} = \mathbb{P}(X=1,Y=4) + \mathbb{P}(X=4,Y=1) = \mathbf{0.12}$
    \item $\mathbf{\mathbb{P}(U=2, V=2)} = \mathbb{P}(X=2,Y=2) = \mathbf{0.02}$
    \item $\mathbf{\mathbb{P}(U=2, V=3)} = \mathbb{P}(X=2,Y=3) + \mathbb{P}(X=3,Y=2) = \mathbf{0.12}$
    \item $\mathbf{\mathbb{P}(U=2, V=4)} = \mathbb{P}(X=2,Y=4) + 0 = \mathbf{0.06}$
    \item $\mathbf{\mathbb{P}(U=3, V=3)} = \mathbb{P}(X=3,Y=3) = \mathbf{0.16}$
    \item $\mathbf{\mathbb{P}(U=3, V=4)} = \mathbb{P}(X=3,Y=4) + 0 = \mathbf{0.12}$
\end{itemize}
\textit{Toutes les autres combinaisons (notamment où $U>V$) ont une probabilité de 0.}

\hrulefill

\section*{Exercice 3 - Lancer de Pièces}

\subsection*{a) Déterminer la loi du couple $(X;Y)$ puis les lois marginales de X et Y.}

\textbf{Raisonnement :} 
L'expérience consiste en 3 lancers consécutifs. L'univers comporte $2^3 = 8$ issues équiprobables (P=Pile, F=Face).
$\Omega = \{PPP, PPF, PFP, PFF, FPP, FPF, FFP, FFF\}$.
$X$ compte les "F" sur les pièces 1 et 2. $Y$ compte les "P" sur les pièces 2 et 3.
En évaluant pour chaque issue (chacune ayant une probabilité de $1/8$) :
\begin{itemize}
    \item $PPP \implies X=0, Y=2$
    \item $PPF \implies X=0, Y=1$
    \item $PFP \implies X=1, Y=1$
    \item $PFF \implies X=1, Y=0$
    \item $FPP \implies X=1, Y=2$
    \item $FPF \implies X=1, Y=1$
    \item $FFP \implies X=2, Y=1$
    \item $FFF \implies X=2, Y=0$
\end{itemize}

\textbf{Reponse :}
On regroupe les probabilités dans un tableau :
\begin{table}[h!]
    \centering
    \begin{tabular}{|c|c|c|c||c|}
    \hline
    $\mathbf{X \setminus Y}$ & \textbf{0} & \textbf{1} & \textbf{2} & \textbf{Loi de X} \\ \hline
    \textbf{0}               & 0          & 1/8        & 1/8        & \textbf{2/8 = 1/4} \\ \hline
    \textbf{1}               & 1/8        & 2/8        & 1/8        & \textbf{4/8 = 1/2} \\ \hline
    \textbf{2}               & 1/8        & 1/8        & 0          & \textbf{2/8 = 1/4} \\ \hline \hline
    \textbf{Loi de Y}        & \textbf{1/4} & \textbf{1/2} & \textbf{1/4} & \textbf{1} \\ \hline
    \end{tabular}
\end{table}

\subsection*{b) Les variables aléatoires sont-elles indépendantes ?}

\textbf{Raisonnement :} 
Vérifions par exemple la case $X=0$ et $Y=0$. D'après le tableau, $\mathbb{P}(X=0, Y=0) = 0$.
Or, le produit des marginales est : $\mathbb{P}(X=0) \times \mathbb{P}(Y=0) = \frac{1}{4} \times \frac{1}{4} = \frac{1}{16}$.

\textbf{Reponse :}
Puisque $0 \neq 1/16$ : \textbf{Les variables X et Y ne sont pas indépendantes.}

\subsection*{c) Calculer $Cov(X;Y)$.}

\textbf{Raisonnement :} 
X suit une loi binomiale sur 2 lancers avec probabilité de succès 1/2. $\mathbb{E}(X) = 1$. De même, $\mathbb{E}(Y) = 1$.
On calcule d'abord $\mathbb{E}(XY)$ en sommant le produit $x \cdot y \cdot p$ pour les cases non nulles du tableau :
$$ \mathbb{E}(XY) = \left(1\cdot 1 \cdot \frac{2}{8}\right) + \left(1\cdot 2 \cdot \frac{1}{8}\right) + \left(2\cdot 1 \cdot \frac{1}{8}\right) = \frac{6}{8} = \frac{3}{4} $$
On utilise ensuite la formule de la covariance.

\textbf{Reponse :}
$$ Cov(X, Y) = \mathbb{E}(XY) - \mathbb{E}(X)\mathbb{E}(Y) = \frac{3}{4} - (1 \times 1) $$
$$ \mathbf{Cov(X, Y) = -\frac{1}{4}} $$

\hrulefill

\section*{Exercice 4 - Lien déterministe non-linéaire}

On considère $X$ telle que $\mathbb{P}(X=-1) = \mathbb{P}(X=0) = \mathbb{P}(X=1) = \frac{1}{3}$. On pose $Y = X^2$.

\subsection*{a) Déterminer la loi de Y.}

\textbf{Raisonnement :} 
Puisque $X \in \{-1, 0, 1\}$, les valeurs possibles pour $Y = X^2$ sont $\{0, 1\}$.
$\mathbb{P}(Y=0) = \mathbb{P}(X=0)$.
$\mathbb{P}(Y=1) = \mathbb{P}(X=-1 \cup X=1) = \frac{1}{3} + \frac{1}{3}$.

\textbf{Reponse :}
\begin{itemize}
    \item $\mathbf{\mathbb{P}(Y=0) = \frac{1}{3}}$
    \item $\mathbf{\mathbb{P}(Y=1) = \frac{2}{3}}$
\end{itemize}

\subsection*{b) Calculer $Cov(X;Y)$.}

\textbf{Raisonnement :} 
\textbf{Etape 1 : Espérances des variables.}\\
$\mathbb{E}(X) = -1\left(\frac{1}{3}\right) + 0 + 1\left(\frac{1}{3}\right) = 0$.
$\mathbb{E}(Y) = 0\left(\frac{1}{3}\right) + 1\left(\frac{2}{3}\right) = \frac{2}{3}$.

\textbf{Etape 2 : Produit des variables.}\\
Le produit $XY$ correspond à $X \cdot X^2 = X^3$.
$\mathbb{E}(XY) = \mathbb{E}(X^3) = (-1)^3\left(\frac{1}{3}\right) + 0^3 + 1^3\left(\frac{1}{3}\right) = 0$.

\textbf{Etape 3 : Formule de la covariance.}\\
$Cov(X,Y) = \mathbb{E}(XY) - \mathbb{E}(X)\mathbb{E}(Y) = 0 - \left(0 \times \frac{2}{3}\right)$.

\textbf{Reponse :}
$$ \mathbf{Cov(X,Y) = 0} $$

\subsection*{c) Les variables X et Y sont-elles indépendantes ?}

\textbf{Raisonnement :} 
Une covariance nulle n'implique pas l'indépendance (la covariance mesure uniquement la corrélation linéaire, or ici la relation est quadratique).
Vérifions avec les probabilités jointes. Regardons l'événement conjoint $\{X=0 \cap Y=1\}$.
Il est impossible d'avoir $X=0$ et $Y=X^2=1$ simultanément. Donc $\mathbb{P}(X=0, Y=1) = 0$.
Cependant, le produit des probabilités marginales donne : $\mathbb{P}(X=0) \times \mathbb{P}(Y=1) = \frac{1}{3} \times \frac{2}{3} = \frac{2}{9}$.

\textbf{Reponse :}
Puisque $0 \neq 2/9$, \textbf{les variables X et Y ne sont pas indépendantes.} (\textit{Leur dépendance est même totale puisque Y est entièrement déterminée par X}).

\end{document}